% Generated draft: requires mathematical and source review.
\documentclass[11pt]{article}
\usepackage[margin=25mm]{geometry}
\usepackage{fontspec}
\setmainfont{Latin Modern Roman}
\usepackage{amsmath,amssymb,mathtools,unicode-math}
\setmathfont{Latin Modern Math}
\usepackage{enumitem,needspace,xurl,hyperref}
\setlength{\parindent}{0pt}
\setlength{\parskip}{5pt}
\setlength{\emergencystretch}{3em}
\newcommand{\N}{\mathbb{N}}
\newcommand{\Z}{\mathbb{Z}}
\newcommand{\Q}{\mathbb{Q}}
\newcommand{\R}{\mathbb{R}}
\newcommand{\C}{\mathbb{C}}
\newcommand{\F}{\mathbb{F}}
\newcommand{\A}{\mathbb{A}}
\newcommand{\PP}{\mathbb{P}}
\newcommand{\E}{\mathbb{E}}
\newcommand{\eps}{\varepsilon}
\newcommand{\dd}{\,\mathrm{d}}
\begin{document}
\section*{UnsolvedMath ID: 800983}
\textbf{Source problem code:} catalogue entry
\subsection{Definitions and mathematical statement}
Let \(M\) be a smooth closed manifold, meaning compact and without boundary. The supplied statement does not specify the signature of \(g\). To interpret its use of light-like vectors and parallel tensors, take \(g\) to be a smooth pseudo-Riemannian metric, with Levi-Civita connection \(\nabla\); any intended restriction on its signature remains unresolved.

A smooth \((1,1)\)-tensor field is a smooth bundle endomorphism \(A:TM\to TM\). It is self-adjoint with respect to \(g\), parallel, and nilpotent if, respectively,
\[
 g(AX,Y)=g(X,AY)\quad\text{for all smooth vector fields }X,Y,
 \qquad \nabla A=0,
 \qquad A^k=0\quad\text{for some integer }k\geq 1.
\]
Assume that such an \(A\) exists with \(A\not\equiv 0\).

\emph{Question (Burns, Matveev, Question 64).} Must there exist either \(N=M\), with \(h=g\), or a smooth two-sheeted covering \(p:N\to M\), with \(h=p^*g\), for which there is a smooth vector field \(V\) satisfying
\[
 V\not\equiv 0,\qquad \nabla^{h}V=0,\qquad h(V,V)=0\quad\text{everywhere on }N?
\]
Here \(\nabla^h\) denotes the Levi-Civita connection of \(h\). The displayed conditions express the requested nonzero light-like parallel vector field using nonzero to mean not identically zero. On a connected manifold a parallel vector field that is not identically zero is nowhere zero. The supplied record does not specify connectedness of the manifold or its double cover, or whether light-like is intended to require a nowhere-zero field on every component; these qualifications require review.
\subsection{Short English statement}
If a closed manifold has a nonzero parallel nilpotent endomorphism of its tangent bundle that is self-adjoint for its metric, must the manifold or a double cover carry a nonzero light-like parallel vector field? The metric signature and connectedness conventions are not specified in the supplied record.
\subsection{Sources}
\begin{itemize}
\item Local UnsolvedMath registry, permanent ID \(800983\), exact source problem code \texttt{catalogue entry}, title: \emph{Open Problems and Questions about Geodesics --- Question 64}. The record attributes the question to Burns and Matveev, with proposed year 2019, and labels it open. That label is supplied registry metadata, not an independently verified mathematical assessment.
\item Burns, Matveev, \emph{Open Problems and Questions about Geodesics} (2019, as identified in the supplied source list), Question 64: \url{https://arxiv.org/abs/1308.5417}. The supplied provenance records access on 2026-10-16 and extraction method \texttt{arxiv\_latex}; these are supplied provenance statements, not a claim of independent source access.
\end{itemize}
\end{document}
